\chapter{Electrical Power System (EPS)}

\section{Requirements and Initial Selection}

The EPS delivers a single regulated 5~V rail to the entire platform from a
two-cell parallel Li-ion pack, with fault protection and battery-voltage
telemetry to the OBC. The design is driven by a small set of requirements:
MR-002 (six-hour continuous flight), PR-008 (power budget margin
$\geq 1.5\times$ against usable, cold-derated capacity), IF-008 (battery
telemetry to the OBC within 5\%), IF-009 (5~V rail, 500~mA continuous, 2~A
peak capability, staying above 4.5~V), SR-001/002/003 (over-current, reverse
polarity, over-discharge protection), MR-007/MR-009 (subsystems within their
own rated ranges; the payload's 0\textdegree{}C floor is scoped to the
payload, not the whole platform), and CR-001 (flyable mass).

\subsection{Battery Selection}

The pack is 2$\times$ Panasonic NCR18650B, unprotected, wired in parallel:
2$\times$3400~mAh at 3.7~V nominal, 25.1~Wh nominal capacity. The 18650
format was chosen for best-in-class energy density with well-characterised
cold performance (rated to $-20$\textdegree{}C, so the battery is not the
platform's thermal bottleneck), and a hermetically sealed steel can that
tolerates the roughly 1~kPa stratospheric environment without the
electrolyte-boiling or pouch-swell failure modes of LiPo. Rejected
alternatives: a USB power bank (poor cold performance, auto-off logic
unsuited to a CubeSat load profile), a LiPo pouch (altitude pressure risk),
and lithium primary cells (non-rechargeable, complicates ground testing).
Because the cells are unprotected, over-current, reverse-polarity, and
over-discharge protection is entirely external --- discrete circuitry plus
firmware (Section~\ref{sec:eps-architecture}).

A boost (not buck-boost) topology is used throughout because the battery
voltage (3.0--4.2~V across full charge/discharge) is always below the 5~V
rail target --- there is no operating point where the input exceeds the
output, so a simpler, single-direction converter suffices.

\subsection{Power Budget: From Estimate to Measurement}

The first budget pass overestimated the payload significantly: it modelled
the Edge TPU running at full inference power (an initially assumed 1.8~W)
for 40\% of the flight. A review corrected two errors: the active inference
power was raised to the datasheet figure of $\sim$3~W (600~mA @ 5~V, higher
than the original estimate), but the inference \emph{duty} was cut from 40\%
to a conservative 5\%, because Edge TPU inference latency for a
MobileNet-class model is on the order of milliseconds against a 1~Hz capture
cycle --- the TPU is only ever briefly active. A second error, a duplicated
boost-efficiency term applied twice in the battery-referred conversion, was
also removed; battery-side draw is now computed with per-rail efficiency
factors applied exactly once ($\sim$1.78$\times$ for 3.3~V loads through the
boost-plus-LDO two-stage conversion, $\sim$1.18$\times$ for 5~V loads through
the single boost stage).

\begin{center}
\begin{tabular}{|l|r|r|}
\hline
\textbf{Quantity} & \textbf{Before review} & \textbf{After review} \\
\hline
Total average battery power & $\sim$2219~mW & $\sim$1216~mW \\
\hline
Energy required (6~h) & 13.3~Wh & $\sim$7.3~Wh \\
\hline
Usable battery (25.1~Wh $\times$ 0.9 cold derate) & 22.6~Wh & 22.6~Wh \\
\hline
Margin & 1.70$\times$ & $\sim$3.1$\times$ \\
\hline
\end{tabular}
\end{center}

The corrected, consolidated budget (referred to the battery, including
conversion losses):

\begin{center}
\begin{tabular}{|l|r|p{7cm}|}
\hline
\textbf{Load} & \textbf{Avg power} & \textbf{Basis} \\
\hline
Coral Dev Board Micro & $\sim$0.76~W & Measured baseline 158--180~mA @ 5~V;
inference adds $\sim$0.3~A for $\sim$2--2.4~ms per cycle, negligible in the
1~Hz average. \\
\hline
Nucleo + sensors + LoRa (3.3~V) & $\sim$0.40~W & Referred through the
boost-plus-LDO two-stage conversion. \\
\hline
Regulator quiescent + series losses & $\sim$0.03~W & Boost quiescent
($\sim$1.5~mA) plus MOSFET/fuse $I^2R$. \\
\hline
\textbf{Total continuous} & \textbf{$\sim$1.2~W} & Consolidated budget,
1195--1216~mW. \\
\hline
\end{tabular}
\end{center}

Against MR-002's six-hour flight this is $\sim$7.2--7.3~Wh required versus
$\sim$22.6~Wh usable (25.1~Wh nominal at a 0.9 cold-derate factor), a margin
of $\sim$3.1$\times$ --- comfortably above the PR-008 floor of 1.5$\times$.
The 0.9 cold-derate factor was, at budget time, an assumption pending
empirical check (Section~\ref{sec:eps-verification}).

\section{Power Architecture}
\label{sec:eps-architecture}

Signal flow: the battery's 2P pack feeds a 4~A polyfuse (FUZETEC
FRG40016F), then an SI2301 P-MOSFET reverse-polarity stage (drain toward the
battery, source toward the load), into a Pololu U3V40F5 5~V boost regulator.
The regulator output feeds the Coral payload, gated by a GPIO-switched
high-side P-MOSFET, and the Nucleo's E5V input; the Nucleo derives 3.3~V
on-board for the sensors, SD card, and LoRa module, so only one regulated
rail exists at the EPS level. A 22~k$\Omega$-top voltage divider with a
100~nF ceramic filter taps the post-fuse battery line to an OBC ADC input
(IF-008); the firmware multiplier is 3.2.

\begin{figure}[htbp]
  \centering
  \includegraphics[width=\textwidth]{eps_schematic.png}
  \caption{EPS and OBC schematic (KiCad).}
\end{figure}

\subsection{Protection Stack}

\begin{itemize}
  \item \textbf{Over-current (SR-001):} a 4~A polyfuse. A boost converter
  output short is fed from the battery through the inductor and output FET
  body diode regardless of converter state, so brief shorts can end before
  the polyfuse's thermal time constant reacts --- observed in practice during
  bench testing.
  \item \textbf{Reverse polarity (SR-002):} the SI2301 stage is the sole
  protection; the U3V40F5 regulator provides none on its own.
  \item \textbf{Over-discharge (SR-003):} layered. Firmware acts first,
  using IF-008 telemetry, at $\sim$3.0~V per cell (orderly shutdown: close
  SD log, send a final telemetry frame, stop inference, halt). A hardware
  EN-pin UVLO is fitted as the backstop below that. The EPS supplies the
  telemetry; FSW owns the shutdown logic.
\end{itemize}

One measured interaction needed a firmware guard: the UVLO recovery point
sits \emph{above} the firmware shutdown threshold, so at end of discharge the
off-load cell voltage relaxes upward, restarts the converter, and reboots
the OBC into a shutdown loop. The mitigation, raised with FSW: at boot, read
the battery voltage first and return to halt if it is still below the
recovery point.

\subsection{Payload Power Control}

The Coral's 5~V feed is gated by a GPIO-controlled high-side P-MOSFET,
giving true power cutoff --- the RT1176 cannot software-cut its own supply,
so the external switch is the only real off. It is paired with a UART
\texttt{SHUTDOWN} command so the payload can quiesce in software first
(close SD logs, send a final frame, drop TPU context) before the rail is
physically cut.

A previously planned bulk capacitor at the Coral connector was deleted after
bench measurement: the real inference transient turned out to be roughly
five times gentler than the worst-case assumption the budget had carried
(Section~\ref{sec:eps-verification}), so the regulator's own onboard
capacitance and control loop absorb the load step without help.

\subsection{Physical Implementation}

The platform stacks as a 10$\times$10$\times$15~cm 1U volume: TT\&C (LoRa)
and CDH occupy the top tier, with batteries, EPS, OBC, and the Coral payload
below. The EPS itself is built on a hand-wired perfboard rather than a PCB
for this iteration, laid out as a single strip carrying the battery input,
polyfuse, reverse-polarity MOSFET, boost regulator, and the battery-telemetry
divider, with dedicated header rows broken out for the Coral power connector
and the OBC power/telemetry connections.

\begin{figure}[htbp]
  \centering
  \includegraphics[width=0.9\textwidth]{eps_perfboard.png}
  \caption{EPS perfboard layout, as built.}
\end{figure}

\section{Test Campaign and Verification}
\label{sec:eps-verification}

The EPS hardware was bench-tested against a 14-step campaign (bring-up
ladder: visual inspection, continuity, current-limited power-up, battery
power-up, then progressively real loads, ending with the endurance run and
an environmental session). The Coral spike test runs first because its
result decides the bulk-capacitor value that gets soldered onto the board;
everything after that tests the board in its near-final configuration. The
bench campaign below is dated 2026-07-07; several tests remain open pending
equipment or OBC integration.

\begin{center}
\small
\begin{tabular}{|l|p{4.3cm}|p{7.7cm}|}
\hline
\textbf{Req.} & \textbf{Requirement} & \textbf{Status} \\
\hline
MR-002 & Operate 6~h or more without external power & Pending: endurance run
not yet executed. \\
\hline
PR-008 & Power budget margin $\geq 1.5\times$ at 6~h, 0.9 cold derate & Shown
by analysis: $\sim$3.1$\times$. Cold-derate assumption still to be checked
empirically (freezer session). \\
\hline
IF-008 & Battery telemetry to OBC within 5\% & Hardware side verified exact
(divider ratio 3.2 confirmed at 3.7~V); ADC-vs-firmware comparison pending. \\
\hline
IF-009 & 5~V rail: 500~mA continuous, 2~A peak, stays above 4.5~V & Preliminary
pass: $\sim$30~mV dip at $\sim$4.91~V during measured inference bursts
(worst-case 0.5~A). Confirmation on the final harness pending. \\
\hline
SR-001 & Over-current protection & Fuse fitted and examined; hold test
pending equipment above 8~A. \\
\hline
SR-002 & Reverse-polarity protection & Verified: reversed input at current
limit draws zero current, rail stays down. \\
\hline
SR-003 & Over-discharge protection, firmware shutdown with hysteresis &
Hardware backstop measured (UVLO 2.8~V cut, 3.1~V recover); firmware test
with FSW pending. \\
\hline
PR-006 & System boot $\leq$30~s (EPS share: rail stable $\leq$2~s) & Pending:
startup capture not yet run. \\
\hline
MR-007 / MR-009 & Subsystems within rated ranges; payload 0\textdegree{}C
floor scoped to payload & Open: freezer and vacuum session scheduled. \\
\hline
MR-008 / CR-001 & Total flyable mass $\leq$650~g & Verified: EPS assembly
15.09~g, well inside its 200~g sub-budget. \\
\hline
\end{tabular}
\end{center}

Headline results from the 2026-07-07 session: the rail regulates at a
constant 5.06~V across the full 3.0--4.2~V battery range (within the
regulator's 4\% accuracy class, on the favourable side); the battery
telemetry divider is exact; the SI2301 reverse-polarity stage fully blocks a
reversed input; and the measured Coral inference transient
($\sim$0.3~A above baseline for 2.0--2.4~ms, self-contained) is far gentler
than the multi-amp spike the original budget assumed, which is what let the
payload-side bulk capacitor be deleted rather than sized up.

\subsection{Open Items}

\begin{itemize}
  \item Confirmation droop capture on the final soldered harness, probing
  directly at the Coral connector, closing the no-added-capacitor decision.
  \item SR-003 firmware shutdown bench test with FSW, including the
  boot-voltage guard that prevents the UVLO-recovery restart loop.
  \item IF-008 ADC-versus-multimeter comparison with the Nucleo at three
  battery voltages.
  \item SR-001 polyfuse hold test now done; trip test needs a source above
  8~A, otherwise falls back to datasheet trip-curve analysis.
  \item PR-006 startup capture, no-load and preloaded.
  \item Six-hour endurance run (MR-002) with a logged discharge curve.
  \item Freezer and vacuum session: functional check plus empirical
  validation of the 0.9 cold-derate assumption.
\end{itemize}

\section{Thermal Risk}

The dominant open risk is thermal, not electrical. The payload's
0\textdegree{}C floor (MR-009) cannot be held through the night-cruise cold
case by passive insulation alone, per thermal analysis --- the working
assumption that $\sim$1.2~W of internal dissipation keeps the interior warm
enough does not survive scrutiny. A small thermostatically controlled
survival heater ($\sim$2--3~W) would close the gap, but its duty cycle draws
directly against the PR-008 margin: at 2.5~W and 30\% duty the total rises to
$\sim$1.95~W and the margin falls to $\sim$1.9$\times$ --- still above the
1.5$\times$ floor, but no longer with the comfortable headroom the 3.1$\times$
figure suggests. Given the limited time remaining before launch, the team
decided to accept the risk of purely passive thermal management rather than
add the heater and its associated budget and integration work. The practical
consequence is layout-driven: every high-internal-heat component (boost
converter, batteries, Coral) is deliberately placed close together so
whatever heat the platform does generate stays concentrated rather than
spreading out to the enclosure walls.